\documentclass[10pt,a4paper]{amsart}
\usepackage[margin=19mm]{geometry}
\usepackage{amsmath,amssymb,mathtools}
\usepackage[scr=rsfs,cal=euler]{mathalfa}
\usepackage{xfrac}
\usepackage[unicode,hidelinks,bookmarks=false]{hyperref}
\newcommand{\ie}{i.e.\ }
\newcommand{\cf}{cf.\ }
\newcommand{\resp}{respectively\ }
\newcommand{\Subsection}{\subsection}
\providecommand{\nobreakdash}{\nobreak\hbox{-}}
\setlength{\emergencystretch}{2em}
\multlinegap0pt
\usepackage{bm}
\usepackage{bbm}
\usepackage{dsfont}
\usepackage{xspace}
\usepackage{cancel}
\usepackage{mathtools}
\usepackage{amsthm}
\makeatletter
\@ifundefined{theorem}{\newtheorem{theorem}{Theorem}}{}
\makeatother
\newcommand{\abs}[1]{\left\vert#1\right\vert}
\newcommand{\cu}{\mathcal U}
\newcommand{\g}{\gamma}
\newcommand{\ps}{\emptyset}
\newcommand{\seq}[1]{\left<#1\right>}
\newcommand{\set}[1]{\left\{#1\right\}}
\newcommand{\sst}{\subseteq}
\renewcommand{\a}{\alpha}
\renewcommand{\b}{\beta}
\renewcommand{\d}{\delta}
\renewcommand{\k}{\kappa}
\renewcommand{\l}{\lambda}
\renewcommand{\o}{\omega}
\title[Topological groups from matrices of sets]{Topological groups from matrices of sets}
\author{Bori\v{s}a Kuzeljevi\'c, Stepan Milo\v{s}evi\'c,, Stevo Todor\v{c}evi\'c}
\date{Source: arXiv:2506.17716v1 (CC BY 4.0)}
\begin{document}
\maketitle

\noindent\emph{Source and licence.} Topological groups from matrices of sets. Version arXiv:2506.17716v1 (CC BY 4.0), distributed under the Creative Commons Attribution 4.0 International licence, as recorded in the arXiv licence metadata for this exact version and shown on the abstract page https://arxiv.org/abs/2506.17716v1. Full rights notice: licenses/arxiv-2506.17716-CC-BY-4.0.txt.\par\smallskip

\noindent\textit{Editorial scope.} Counted unit: the Theorem whose source label is \texttt{t:frechet} in the article (source0.tex lines 442--445, source label \texttt{t:frechet}), delivered together with the article's own complete original proof of it (source0.tex lines 446--519, one \texttt{proof} environment of 74 lines). No mathematics is written, repaired or re-derived: statement and proof are the article's own text line for line. Required context retained uncounted: t:osnovna (lines 232--243); t:smallcharacter (lines 376--379). Every definition, display and label of those lines is retained unchanged. Omitted from this package and not redistributed: 1--231, 244--375, 380--441, 520--1211. Nothing inside a retained excerpt is omitted or altered: every retained body is delivered exactly as the article's lines are written, so no omission receipt of any kind accompanies this package. Citations: the retained excerpts carry no citation command at all, so no bibliography entry of the article is transcribed and no reference is redistributed. Reference adaptation: none was needed and none was made; every reference inside the delivered unit resolves inside it, so no unresolved reference or citation is produced. Printed slips kept literal: no printed word, symbol, hypothesis or number of the counted statement or of its proof was altered. Layout and class: the article's own class is \texttt{amsart} with packages including amssymb, amsthm, amsmath, hyperref; the wrapper is the lane's pinned \texttt{amsart} editorial wrapper at 10pt on A4, which restates the theorem-like environments the retained text needs and adds the article's own macro definitions that the retained text needs (\texttt{\textbackslash a}, \texttt{\textbackslash abs}, \texttt{\textbackslash b}, \texttt{\textbackslash cu}, \texttt{\textbackslash d}, \texttt{\textbackslash g}, \texttt{\textbackslash k}, \texttt{\textbackslash l}, \texttt{\textbackslash o}, \texttt{\textbackslash ps}, \texttt{\textbackslash seq}, \texttt{\textbackslash set}, \texttt{\textbackslash sst}), with extra wrapper packages (bm, bbm, dsfont, xspace, cancel, mathtools). The delivered numbering is the unit's own. Assets: no retained span contains a figure environment, a graphics-inclusion command, a PDF-page inclusion, a table or a TikZ picture; no image file is copied into this package, the package declares no asset dependency and no image file is redistributed with it. Licence: version arXiv:2506.17716v1 of this article is distributed under the Creative Commons Attribution 4.0 International licence, as recorded in the arXiv licence metadata for this exact version and shown on the abstract page https://arxiv.org/abs/2506.17716v1. A full rights notice is delivered at licenses/arxiv-2506.17716-CC-BY-4.0.txt. Characters: every retained excerpt is pure ASCII, so the delivered engine, XeLaTeX with its default Latin Modern text font, reports no missing character.
\begin{theorem}\label{t:osnovna}
	Let $G$ be a group with the identity $e$, and $\cu$ a family of subsets of $G$ satisfying the following conditions:
	\begin{enumerate}
		\item for every $U\in\cu$ there is an element $V\in\cu$ such that $V^2\sst U$;
		\item for every $U\in\cu$ there is an element $V\in\cu$ such that $V^{-1}\sst U$;
		\item for every $U\in\cu$ and every $x\in U$, there is $V\in\cu$ such that $Vx\sst U$;
		\item for every $U\in\cu$ and $x\in G$, there is $V\in\cu$ such that $xVx^{-1}\sst U$;
		\item for $U,V\in\cu$, there is $W\in\cu$ such that $W\sst U\cap V$;
		\item $\set{e}=\bigcap \cu$.
	\end{enumerate}
	Then $G$ is a topological group with the topology generated by the base $\set{Ua:a\in G}$, i.e. $\cu$ is an open base of the identity in $G$.
\end{theorem}

\begin{theorem}\label{t:smallcharacter}
    Suppose $\l\le\k$ are infinite regular cardinals and $\mathcal F$ is a strong $(\k,\l)$-matrix.
	Then for $\d\in C$, the family $\cu_{\d}(\mathcal F)$ is a local base of $\ps$ in the group $([\d]^{<\o}, \triangle)$ with the topology induced from $G(\mathcal F)$.
\end{theorem}

\begin{theorem}\label{t:frechet}
	If $\l\le\k$ are regular infinite cardinals, and $\mathcal F$ is a strong $(\k,\l)$-matrix indexed by $\k^+$, then 
	$G(\mathcal F)$ is $\l$-Frech\'et if and only if $\mathcal F$ is unbounded.
\end{theorem}

\begin{proof}
	By Theorem \ref{t:osnovna}, we know that $G(\mathcal F)$ is a topological group.
	We first assume that the matrix $\mathcal F$ is unbounded.
	To show that $G(\mathcal F)$ is $\l$-Frech\'et, take $A\sst G(\mathcal F)$ such that $\ps$ is in the closure of $A$.
	Let $\theta$ be a regular cardinal above $\left(2^{\o_1}\right)^+$.
	Take $M$, an elementary submodel of $(H(\theta),\in)$ of size $\k$, containing $\mathcal F$ and $A$ as elements.
	Denote $\d=M\cap \k^+$.
	Note that $\d$ is an ordinal in $\k^+$ and that $\l\sst\d\sst M$.
	To prove this direction of the theorem, it is enough to show that $\ps$ is in the closure of $M\cap A$.
	To see this, note that by Theorem \ref{t:smallcharacter}, we know that for every basic open neighborhood of $\ps$ in $G(\mathcal F)$, there is some $\xi<\l$ such that the induced open set on $[\d]^{<\o}$ is exactly $[\d]^{<\o}\cap U_{\xi}(\d)$.
	If $\ps$ were in the closure of $A\cap M$, we would be able to pick $a_{\xi}\in U_{\xi}(\d)\cap A\cap M$ for each $\xi<\l$.
	Then the $\l$-sequence $\seq{a_{\xi}:\xi<\l}$ would be in $A$.
	Moreover, for any open $U_{\eta}(\d)$ and each $\xi>\eta$ we would have $a_{\xi}\in U_{\xi}(\d)\sst U_{\eta}(\d)$ (this holds as $F_{\eta}(\d)\sst F_{\xi}(\d)$ by the definition of matrix).
	Thus it would prove that $G(\mathcal F)$ is $\l$-Frech\'et.

	So we continue to prove $\ps\in \overline{A\cap M}$.
	If $\abs{A}<\k^+$, then $A\sst M$ and by the above observation the theorem would be proved, so assume that $A$ is of cardinality $\k^+$.
	We now work towards proving $\ps\in\overline{A\cap M}$.
	So take $\xi<\l$.
	Note that $\xi\in M$.
	We will prove that $U_{\xi}(\d)\cap A\cap M\neq\ps$, i.e. that there is some $a\in A\cap M$ such that $a\cap F_{\xi}(\d)=\ps$.
	Since $\ps\in \overline{A}$, there is some $b\in A$ such that $b\cap F_{\xi}(\d)=\ps$.
	We can assume the $b\notin M$ (otherwise the proof is finished).
	Denote $c=b\setminus M$.
	
	Consider the formula $$\psi(\a)\equiv \left(\exists w\in [\k^+]^{<\o}\right)\ \a<w\ \wedge\ w\cup (b\cap \d)\in A.$$
	%Let us consider the family $$\mathcal A_{b\cap\d}=\set{w\in [\k^+]^{<\o}:w\cup (b\cap \d)\in A}.$$
	To see that $M\models \big((\forall \a<\k^+)\psi(\a)\big)^M$, take arbitrary $\a\in M\cap \k^+=\d$ for a moment.
	Then $H(\theta)\models \psi(\a)$, namely $b\setminus \d$ is a witness for this.
	Since $A,[\k^+]^{<\o},b\cap\d$, and $\a$ are all elements of $M$, by elementarity, $\psi(\a)$ holds in $M$.
	Since $\a$ was arbitrary in $M\cap \k^+$, we proved $M\models \big((\forall \a<\k^+)\psi(\a)\big)^M$.
	Again, by elementarity, $H(\theta)\models (\forall \a<\k^+)\psi(\a)$.
	This means that, in $H(\theta)$, for each $\a<\k^+$ there is some $w_{\a}$ such that $\a\cap w_{\a}=\ps$ and $w_{\a}\cup (b\cap\d)\in A$.
	Moreover, we can choose $w_{\a}$ to be minimal with respect to the anti-lexicographical ordering of $[\k^+]^{<\o}$.
	As this ordering is an element of $M$, it means that the function $f:\k^+\to [\k^+]^{<\o}$ given by $f(\a)=w_{\a}$ is definable in $M$ by parameters in $M$, so $f\in M$.
	Hence, the family $\mathcal A=\set{w_{\a}:\a<\k^+}$ is an element of $M$.

	To finish the proof, since $b\cap F_{\xi}(\d)=\ps$, we need to show that there is some $\a<\d$ such that $w_{\a}\cap F_{\xi}(\d)=\ps$.
	Equivalently, we need $\a<\d$ such that $\rho_{\mathcal F}(x,\d)>\xi$ for each $x\in w_{\a}$.
	Suppose this is not the case, i.e. $$(\forall \a<\d)(\exists x\in w_{\a})\ \rho_{\mathcal F}(x,\d)\le\xi.$$
	%Fix for a moment $\a<\d$.
	Consider the family $\mathcal C=\set{\b<\k^+:(\forall \a<\b)(\exists x\in w_{\a})\rho_{\mathcal F}(x,\b)\le\xi}$.
	Since $\mathcal F\in M$, we know that $\rho_{\mathcal F}\in M$.
	So, as $\xi\in M$ as well, we have $\mathcal C\in M$.
	From $\d\in \mathcal C\setminus M$ (by the assumption), we know that $\mathcal C$ is cofinal in $\k^+$. %,thus $M\cap \mathcal C$ is cofinal in $\d=\k^+\cap M$.
	%Thus, in $M$, for each $\a\in \d$ there is $\g\ge\a$ in $M$ such that $\rho_{\mathcal F}(x,\g)\le \xi$ for some $x\in w_{\a}$.
	Hence, for each $\a<\k^+$ there is some minimal $\g_{\a}\ge\a$ such that $\rho_{\mathcal F}(x,\g_{\a})\le \xi$ for some $x\in w_{\a}$.
	%So, in $H(\theta)$, for each $\a<\k^+$ take minimal $\g_{\a}\ge\a$ such that $\rho_{\mathcal F}(x,\g_{\a})\le \xi$.
	Now consider the family $\mathcal A'=\set{w_{\a}\cup\set{\g_{\a}}:\a<\k^+}$.
	This is an uncountable family of pairwise disjoint finite subsets of $\k^+$.
	%We can assume that the size of all finite sets in $\mathcal A'$ is the same.
	Since $\rho_{\mathcal F}$ is $\k^+$-unbounded (the assumption of the theorem), there are $\a<\b<\k^+$ such that $\rho_{\mathcal F}(x,y)>\xi$ for each $x\in w_{\a}\cup\set{\g_{\a}}$ and $y\in w_{\b}\cup\set{\g_{\b}}$.
	But, this in particular means that there is no $x\in w_{\a}$ such that $\rho_{\mathcal F}(x,\g_{\b})\le \xi$, contradicting the choice of $\g_{\b}$.
	Hence, there is some $\a_0<\d$ such that $\rho_{\mathcal F}(x,\d)>\xi$ for each $x\in w_{\a_0}$.
	Then, since $\a_0<\d$, we know that $w_{\a_0}\in M\cap \mathcal A$ and $w_{\a_0}\cap F_{\xi}(\d)=\ps$.
	Consequently $(b\cap\d)\cup w_{\a_0}\in A\cap M$ and $\left((b\cap\d)\cup w_{\a_0}\right)\cap F_{\xi}(\d)=\ps$, as required.
	So, the proof of one direction of the theorem is completed.

	For the other direction, let us assume that $G(\mathcal F)$ is $\l$-Frech\'et.
	Suppose that $\rho_{\mathcal F}$ is not $\k^+$-unbounded.
	This means that there are $\xi_0<\l$ and a family $\mathcal A$ of size $\k^+$ consisting of pairwise disjoint finite subsets of $\k^+$ such that for all different $a$ and $b$ in $\mathcal A$ there are $\a\in a$ and $\b\in b$ such that $\rho_{\mathcal F}(\a,\b)\le \xi_0$.
	Then $\mathcal A\sst G(\mathcal F)$.
	Let us show that $\ps\in\overline{\mathcal A}$.
	Take any basic open $U_{\eta}(\b)$.
	Since $\mathcal A$ is a family of pairwise disjoint finite sets of size $\k^+$, there is some $a\in \mathcal A$ with $\b<a$.
	Then $a\cap F_{\eta}(\b)=\ps$, i.e. $a\in \mathcal A\cap U_{\eta}(\b)$.
	So, since $G(\mathcal F)$ is $\l$-Frech\'et, there is a $\l$-sequence $\seq{a_{\xi}:\xi<\l}$ in $\mathcal A$ such that: \begin{equation}\label{eq:convergence}\mbox{for any open }U_{\zeta}(\g)\mbox{ there is }\eta<\l\mbox{ so that }a_{\xi}\in U_{\zeta}(\g)\mbox{ whenever }\eta\le\xi<\l.\end{equation}
	Take $a\in \mathcal A$ such that $a_{\xi}<a$ for each $\xi<\l$.
	This is possible as $\mathcal A$ is of size $\k^+$ and $\l\le\k$.
	Enumerate $a=\set{\b_0,\dots,\b_{l-1}}$ for some $l<\o$.
	By the assumption on $\mathcal A$, for each $\xi<\l$ there is some $i<l$ and $\a\in a_{\xi}$ such that $\rho_{\mathcal F}(\a,\b_i)\le \xi_0$.
	Hence, there is $j<l$ and a set $J$ cofinal in $\l$, so that for $\xi\in J$ there is $\a\in a_{\xi}$ such that $\rho_{\mathcal F}(\a,\b_j)\le \xi_0$.
	But this implies that $a_{\xi}\notin U_{\xi_0}(\b_j)$ for each $\xi\in J$, contradicting the assumption that $\seq{a_{\xi}:\xi<\l}$ satisfies (\ref{eq:convergence}) for $\zeta=\xi_0$ and $\g=\b_j$.
\end{proof}

\end{document}
